\section{Studi Kasus: Basis Pengetahuan Silsilah Keluarga}

Studi kasus ini menyajikan pemodelan silsilah keluarga menggunakan logika predikat orde pertama. Kasus ini dipilih karena mengilustrasikan secara jelas bagaimana entitas, properti, relasi, dan aturan saling berinteraksi dalam basis pengetahuan \cite{rosen2019}.

\subsection{Kebutuhan pemodelan}

Basis pengetahuan silsilah dimodelkan dengan relasi-relasi berikut:

\begin{itemize}[leftmargin=*, nosep]
    \item \textbf{Fakta dasar} $OrangTua(x, y)$: entitas $x$ adalah orang tua dari $y$.
    \item \textbf{Fakta sifat} $Laki(x)$ dan $Perempuan(x)$: jenis kelamin entitas $x$.
    \item \textbf{Aturan} $Ayah(x, y)$: $x$ ayah dari $y$ jika $x$ orang tua dari $y$ dan $x$ laki-laki.
    \item \textbf{Aturan} $Ibu(x, y)$: $x$ ibu dari $y$ jika $x$ orang tua dari $y$ dan $x$ perempuan.
    \item \textbf{Aturan} $Kakek(x, z)$: $x$ kakek dari $z$ jika $x$ ayah dari suatu $y$ dan $y$ orang tua dari $z$.
    \item \textbf{Aturan} $Saudara(x, y)$: $x$ saudara $y$ jika ada $z$ yang merupakan orang tua keduanya dan $x \neq y$.
\end{itemize}

\subsection{Deklarasi Formal}

Himpunan fakta ($\Delta$):
\begin{align*}
\Delta = \{ & OrangTua(Hasan, Ahmad), \; OrangTua(Hasan, Rina), \\
            & OrangTua(Minah, Ahmad), \; OrangTua(Minah, Rina), \\
            & OrangTua(Ahmad, Ali), \; OrangTua(Ahmad, Sari), \\
            & Laki(Hasan), \; Laki(Ahmad), \; Laki(Ali), \\
            & Perempuan(Minah), \; Perempuan(Rina), \; Perempuan(Sari) \}
\end{align*}

Himpunan aturan ($\Gamma$):
\begin{align*}
R_1&: \forall x, y \; (OrangTua(x, y) \wedge Laki(x) \rightarrow Ayah(x,y)) \\
R_2&: \forall x, y \; (OrangTua(x, y) \wedge Perempuan(x) \rightarrow Ibu(x,y)) \\
R_3&: \forall x, y, z \; (Ayah(x, y) \wedge OrangTua(y, z) \rightarrow Kakek(x,z)) \\
R_4&: \forall x, y, z \; (OrangTua(z, x) \wedge OrangTua(z, y) \wedge (x \neq y) \rightarrow Saudara(x, y))
\end{align*}

\subsection{Resolusi Query: Siapa Kakek dari Ali?}

\begin{contoh}
\textbf{Query}: $Kakek(?, Ali)$ --- siapa yang merupakan kakek dari Ali?

\textbf{Penalaran (Backward Chaining dari $R_3$):}
\begin{enumerate}
    \item Aturan $R_3$ menyatakan: $Kakek(x, z)$ jika $Ayah(x, y) \wedge OrangTua(y, z)$.
    \item Untuk $z = Ali$, cari $y$ sehingga $OrangTua(y, Ali)$. Dari $\Delta$: $y = Ahmad$.
    \item Cari $x$ sehingga $Ayah(x, Ahmad)$. Gunakan $R_1$: perlu $OrangTua(x, Ahmad) \wedge Laki(x)$.
    \item Dari $\Delta$: $OrangTua(Hasan, Ahmad)$ dan $Laki(Hasan)$. Maka $Ayah(Hasan, Ahmad)$.
    \item Substitusi ke $R_3$: $Ayah(Hasan, Ahmad) \wedge OrangTua(Ahmad, Ali) \rightarrow Kakek(Hasan, Ali)$.
    \item Anteseden terpenuhi, maka $Kakek(Hasan, Ali)$ bernilai \textbf{Benar}. $\blacksquare$
\end{enumerate}
\end{contoh}

\subsection{Resolusi Query: Siapa Saudara dari Ahmad?}

\begin{contoh}
\textbf{Query}: $Saudara(?, Ahmad)$

\textbf{Penalaran dari $R_4$:}
\begin{enumerate}
    \item $R_4$: $Saudara(x, y)$ jika $OrangTua(z, x) \wedge OrangTua(z, y) \wedge (x \neq y)$.
    \item Untuk $y = Ahmad$, cari $z$ sehingga $OrangTua(z, Ahmad)$. Dari $\Delta$: $z = Hasan$ atau $z = Minah$.
    \item Dengan $z = Hasan$: cari $x \neq Ahmad$ sehingga $OrangTua(Hasan, x)$. Dari $\Delta$: $x = Rina$.
    \item Verifikasi: $OrangTua(Hasan, Rina) \wedge OrangTua(Hasan, Ahmad) \wedge (Rina \neq Ahmad)$. Semua terpenuhi.
    \item Kesimpulan: $Saudara(Rina, Ahmad)$ bernilai \textbf{Benar}. $\blacksquare$
\end{enumerate}
\end{contoh}

\begin{konsep}
Studi kasus ini menunjukkan bahwa entitas (Hasan, Ahmad, dan lain-lain) dimodelkan sebagai konstanta, relasi (\textit{OrangTua}) sebagai predikat arity-2, sifat (\textit{Laki}, \textit{Perempuan}) sebagai predikat arity-1, dan aturan dinyatakan sebagai implikasi berkuantor universal ($\forall$). Pola penalaran demikian menjadi dasar sistem pakar dan basis data deduktif.
\end{konsep}
